\chapter{Transformada de Fourier de tiempo discreto}

\section{Introducci\'{o}n}

\begin{enumerate}

\item   Vamos a ver como podemos \textbf{modificar} el concepto de \textbf{serie de Fourier} que 
relacionamos normalmente con las funciones peri\'{o}dicas de tiempo discreto 
\textbf{para estudiar funciones no peri\'{o}dicas} de tiempo discreto. 
Esto nos permite \textbf{definir} el concepto de \textbf{contenido de frecuencias de una se\~{n}al}, 
a\'{u}n cuando sea aperi\'{o}dica.

\item   Pensemos en una se\~{n}al $\hat{x}\left[ n\right]$ peri\'{o}dica de tiempo discreto con un
        per\'{\i}odo fundamental $2 N + 1$ definida por
        \begin{equation}
                        \begin{array}{ccc}
                                -N \leq n < -N_{1} &\rightarrow& \hat{x}[n] = 0, \\
                                -N_{1} \leq n \leq N_{1} & \rightarrow& \hat{x}[n] = 1, \\
                                N_{1} < n \leq N &\rightarrow& \hat{x}[n] = 0.
                        \end{array}
        \end{equation}
        \begin{equation}
                \hat{x} \left[ n\right] = \hat{x} \left[n + 2 \, N + 1\right].
        \end{equation}

\item   Recordemos la ecuaci\'{o}n de an\'{a}lisis de la serie de Fourier de tiempo discreto: 
  para este caso podemos escribir as\'{\i} 
        \begin{equation}
                a_{k}=
                        \frac{1}{2 \, N + 1} \, \sum_{n=-N_{1}}^{N_{1}} \hat{x}\left[ n\right] \, e^{-j \, k \, \omega_{N} \, n},
        \end{equation}
donde $\omega_{N} = \frac{2 \, \pi}{2 \, N + 1}$.
        
\item   Para el caso $k=0$ obtenemos el coeficiente $a_{0}$ (com\'{u}nmente llamado en 
  electr\'{o}nica nivel de continua)
  \begin{equation}
    a_{0} = \frac{1}{2 \, N + 1} \sum_{n=-N_{1}}^{N_{1}} 1 = \frac{2 \, N_{1} + 1}{2 \, N + 1}.
  \end{equation}
  El valor  $\frac{2 \, N_{1} + 1}{2 \, N + 1}$ es denominado ciclo de trabajo del tren de pulsos, e indica
el porcentaje del per\'{\i}odo en que la se\~{n}al es igual a $1$.
        
\item   Para el caso $k \neq 0$ obtenemos los dem\'{a}s coeficientes $a_{k}$
  \begin{equation}
    a_{k}=  
    \frac{1}{2 \, N + 1} \sum_{n=-N_{1}}^{N_{1}} e^{-j \, k \, \omega_{N} \, n}.
  \end{equation}
  
\item   Multiplicando ambos lados de la ecuaci\'{o}n por $2 \, N + 1$ y desdoblando la sumatoria resulta
  \begin{equation}
    (2 \, N + 1) \, a_{k} =
    \sum_{n=-N_{1}}^{-1} e^{-j \, k \, \omega_{N} \, n} +
    \sum_{n=1}^{N_{1}} e^{-j \, k \, \omega_{N} \, n} + 1,
  \end{equation}
  que podemos escribir como
  \begin{equation}
    (2 \, N + 1) \, a_{k} =
    \sum_{n=-N_{1}}^{0} e^{-j \, k \, \omega_{N} \, n} +
    \sum_{n=0}^{N_{1}} e^{-j \, k \, \omega_{N} \, n} - 1,
  \end{equation}
  teniendo en cuenta  que al poner un l\'{\i}mite igual a $0$ en ambas 
  sumatorias  estamos sumando dos veces un t\'{e}rmino $e^{0} = 1$ y necesitamos
  $2$ para no alterar la igualdad.

\item   Invertimos los l\'{\i}mites en la primer sumatoria e invertimos el signo del exponente 
  correspondiente para no alterar el resultado
  \begin{equation}
    (2 \, N + 1) \, a_{k} =
    \sum_{n=0}^{N_{1}} e^{j \, k \, \omega_{N} \, n} +
    \sum_{n=0}^{N_{1}} e^{-j \, k \, \omega_{N} \, n} - 1.
  \end{equation}
        
\item   Tenemos entonces dos sumatorias de sucesiones geom\'{e}tricas. Por lo  tanto podemos  escribir
  \begin{equation}
    (2 \, N + 1) \, a_{k} =
    \frac{1 - e^{j \, k \, \omega_{N} \, (N_{1}+1)}}{1 - e^{j \, k \, \omega_{N}}}  +
    \frac{1 - e^{-j \, k \, \omega_{N} \, (N_{1}+1)}}{1 - e^{-j \, k \, \omega_{N}}}  -
    1. \label{ejemploTFTDinicial}
  \end{equation}
  
\item  La utilidad de la ecuaci\'{o}n (\ref{ejemploTFTDinicial}) se ve si la escribimos as\'{\i}
  \begin{equation}
    (2 \, N + 1) \, a_{k} = \left. X(e^{j \omega}) \right|_{\omega = k \, \omega_{N}} =
    \left.
      \frac{1 - e^{j \, \omega \, (N_{1}+1)}}{1 - e^{j \, \omega}}  +
      \frac{1 - e^{-j \, \omega \, (N_{1}+1)}}{1 - e^{-j \, \omega}}  -
      1
    \right|_{\omega=k \, \omega_{N}},  \label{funTFTDinicial}
  \end{equation}
  y entonces es evidente que la podemos interpretar como una \textbf{secuencia} formada por
  \textbf{muestras} de una  \textbf{funci\'{o}n envolvente}. La variable $\omega $ es continua y la
  funci\'{o}n continua $X(e^{j \, \omega})$ muestreada en (\ref{funTFTDinicial}) es la envolvente de la secuencia 
  $N a_{k}$, formada por valores igualmente espaciados en frecuencia.
  
\item   Observe que la separaci\'{o}n de las muestras de la envolvente de $N a_{k}$ est\'{a} dada por
  \begin{equation}
    \omega_{N} = \frac{2 \, \pi}{2 \, N + 1}.
  \end{equation}
  Entonces, cuando mayor sea $N$ menor ser\'{a} $\omega_{N}$ y la cantidad de muestras en un 
  intervalo dado de frecuencia ser\'{a} mayor.
  
\item  Si fijamos el valor $N_{1}$ (recordando que debe ser $N_{1} < N,$ por hip\'{o}tesis) la 
  envolvente de $N \, a_{k}$ es independiente de $N$.
        
\end{enumerate}

\section{Las ecuaciones de an\'{a}lisis y s\'{\i}ntesis}

\subsection{Análisis}

\begin{enumerate}

\item   Definamos una se\~{n}al $x\left[ n\right] $ no peri\'{o}dica  
\begin{eqnarray}
  \left( 0\leq n\leq N_{1}\right)                         &\rightarrow    & x\left[ n\right] \neq 0, \\
  \left( n < 0\right) \vee \left( N_{1} < n\right)            &\rightarrow    & x\left[ n\right] =0.
\end{eqnarray}

\item   Definamos una se\~{n}al peri\'{o}dica $\tilde{x}\left[ n\right] =\tilde{x}\left[ n+N\right]$ 
(donde $N > N_{1}$) de forma tal que 
\begin{eqnarray}
 \left( n \bmod N\right) > N_{1}         &\rightarrow& \tilde{x}\left[ n\right] =0, \\ 
 \left( n \bmod N\right) \leq N_{1}      &\rightarrow& \tilde{x}\left[ n\right] = x \left[ n\bmod N\right].
\end{eqnarray}

\item   La se\~{n}al peri\'{o}dica $\tilde{x}\left[ n\right]$ se puede escribir como una serie de Fourier 
\begin{equation}
        \tilde{x}\left[ n\right] =\sum_{k=0}^{N-1}a_{k} \, e^{j \, k \, \left( \frac{2 \, \pi}{N}\right) \, n},
\end{equation}
donde 
\begin{equation}
  N \, a_{k} = \sum_{n=0}^{N-1}\tilde{x}\left[ n\right] \, e^{-j \, k \, \left( \frac{2 \, \pi}{N}\right) \, n}.
\end{equation}

\item Entonces como $\tilde{x}\left[ n\right] =x\left[ n\right] $ en el intervalo $[0,N-1]$ se puede 
reemplazar $\tilde{x}\left[ n\right] $ por $x\left[ n\right]$ 
\begin{equation}
  N \, a_{k}=
  \sum_{n=0}^{N-1}\tilde{x}\left[ n\right] \, e^{-j \, k \, \left( \frac{2 \, \pi }{N}\right) \, n}
  =
  \sum_{n=0}^{N-1}x\left[ n\right] e^{-j \, k \, \left( \frac{2 \, \pi}{N}\right) \, n}.
\end{equation}

\item Adem\'{a}s como $x\left[ n\right] =0$ fuera del intervalo $\left[ 0,N_{1}\right]$ podemos cambiar 
los l\'{\i}mites de la sumatoria 
\begin{equation}
  N \, a_{k}=
  \sum_{n=0}^{N-1}x\left[ n\right] \, e^{-j \, k \, \left( \frac{2 \, \pi}{N}\right) \, n}=
  \sum_{n=-\infty }^{\infty }x\left[ n\right] \, e^{-j \, k \, \left( \frac{2 \, \pi }{N}\right) \, n}.
\end{equation}

\item   Definimos la ecuaci\'{o}n de an\'{a}lisis de la transformada de Fourier de tiempo discreto
\begin{equation}
    X\left( e^{j \, \omega }\right) = \sum_{n=-\infty }^{\infty }x\left[ n\right] \, \left(e^{j \, \omega}\right)^{-n}.
\end{equation}

\end{enumerate}

\subsection{S\'{\i}ntesis}

\begin{enumerate}
\item Entonces tenemos que los coeficientes de Fourier se pueden calcular
\begin{equation}
  N \, a_{k}= X\left( e^{j \, k \, \omega_{N}}\right),
\end{equation}
donde 
\begin{equation}
 \omega_{N} = \frac{2 \, \pi}{N}.
\end{equation}

\item   Por lo tanto podemos reescribir la serie de Fourier como
\begin{equation}
  \tilde{x}\left[ n\right] =
  \sum_{k=0}^{N-1}\frac{1}{N} \, X\left(e^{j \, k \, \omega_{N}}\right) \, e^{j \, k \, \omega_{N} \, n}.
\end{equation}
Dado que $\omega_{N} = \frac{2 \, \pi}{N}$ tenemos que $\frac{\omega_{N}}{2 \, \pi}=\frac{1}{N}$ y por lo tanto 
\begin{equation}
  \tilde{x}\left[ n\right] =
  \frac{1}{2 \, \pi} \, \sum_{k=0}^{N-1}X\left(e^{j \, k \, \omega_{N}}\right) \, e^{j \, k \, \omega_{N} \, n} \, \omega_{N}.
\end{equation}

A medida que la cantidad de puntos $N\rightarrow \infty$, 
la frecuencia tiende a un diferencial $\omega_{N}\rightarrow \, d\omega$, 
la funci\'{o}n peri\'{o}dica se parece cada vez m\'{a}s a la no peri\'{o}dica
$\tilde{x}\left[ n\right] \rightarrow x\left[ n\right],$ y 
se obtiene la ecuaci\'{o}n de s\'{\i}ntesis 
\begin{equation}
  x\left[ n\right] =
  \frac{1}{2 \, \pi } \, \int_{0}^{2 \, \pi } X\left( e^{j \, \omega }\right) \, e^{j \, \omega \, n} \, \, d\omega .
\end{equation}
\end{enumerate}

\section{Convergencia de la Transformada de Fourier}

La transformada de Fourier 
\begin{equation}
  X\left( e^{j \, \omega }\right) =
  \sum_{n=-\infty }^{\infty }x\left[ n\right]e^{-j \, \omega \, n},
\end{equation}
de una secuencia $x\left[ n\right] $ existe si se cumple 
\begin{equation}
\left( \sum_{n=-\infty }^{\infty }\left| x\left[ n\right] \right| \right)
<M_{1},
\end{equation}
\textbf{o} si se cumple 
\begin{equation}
\left( \sum_{n=-\infty }^{\infty }\left( x\left[ n\right] \right)^{2}\right) <M_{2},
\end{equation}
donde $M_{1}$ y $M_{2}$ son dos n\'{u}meros reales positivos finitos.

Observe que la ecuaci\'{o}n de s\'{\i}ntesis no presenta problemas de
convergencia pues sus l\'{\i}mites de integraci\'{o}n est\'{a}n restringidos
a un per\'{\i}odo 
\begin{equation}
x\left[ n\right] =
  \frac{1}{2 \, \pi }\int_{0}^{2 \, \pi }
  X\left( e^{jk\omega}\right) e^{jk\omega \, n}\, d\omega .
\end{equation}

\section{Transformada de Fourier para se\~{n}ales peri\'{o}dicas}

\begin{enumerate}

\item Consideremos la se\~{n}al 
\begin{equation}
  x\left[ n\right] =e^{j \, \omega_{0} \, n}.
\end{equation}

\item Podemos probar que su transformada de Fourier es 
\begin{equation}
  X\left( e^{j \, \omega }\right) =
  2 \, \pi \, \left( \sum_{k=-\infty }^{\infty }
    \delta \left(\omega -\omega _{0}-2 \, \pi \, k\right) \right).
\end{equation}

\item Para hacerlo, usemos la ecuaci\'{o}n de s\'{\i}ntesis 
\begin{equation}
\frac{1}{2 \, \pi }\int_{-\pi}^{\pi }
  X\left( e^{j \, \omega }\right) e^{j \, \omega \, n}\, d\omega =
  \frac{2 \, \pi}{2 \, \pi }\int_{-\pi}^{\pi}
    \left(
      \sum_{k=-\infty}^{\infty } \delta \left( \omega -\omega _{0}-2 \, \pi \, k \right) 
    \right)
    e^{j \, \omega \, n}\, d\omega .
\end{equation}

\item Observemos que el intervalo $\left[-\pi, \pi\right] $ contiene exactamente un
impulso en $\omega = \omega_{0} + 2 \, \pi \, k$, por lo tanto 
\begin{equation}
\frac{1}{2 \, \pi }\int_{-\pi}^{\pi }X\left( e^{j \, \omega }\right) e^{j \, \omega \, n}\, d\omega =
  \int_{-\pi}^{\pi }\delta \left( \omega -\omega _{0} - 2 \, \pi \, k\right)e^{j \, \omega \, n}\, d\omega =
  e^{j \, \left(\omega_{0} + 2 \, \pi \, k\right) \, n}.
\end{equation}

\item En general, a la se\~{n}al peri\'{o}dica 
\begin{equation}
x\left[ n\right] =\sum_{k=0}^{N-1}a_{k} \, e^{j \, k\left( \frac{2 \, \pi }{N}\right) \, n},
\end{equation}
le corresponde la transformada
\begin{equation}
X\left( e^{j \, \omega }\right) =
2 \, \pi \, \left(\sum_{k=-\infty }^{\infty } a_{k} \, \delta\left( \omega -\frac{2 \, \pi \, k}{N}\right)\right). 
\label{tfdperiodica}
\end{equation}

\item Para probarlo, usemos la ecuaci\'{o}n de s\'{\i}ntesis 
\begin{equation}
\frac{1}{2 \, \pi }\int_{0}^{2 \, \pi}X\left( e^{j \, \omega }\right) \, e^{j \, \omega \, n}\, d\omega =
\frac{1}{2 \, \pi }\int_{0}^{2 \, \pi} 
  \left( 
    \sum_{k=-\infty }^{\infty} 2 \, \pi \, a_{k} \, \delta \left( \omega -\frac{2 \, \pi \, k}{N} \right) 
  \right) \, e^{j \, \omega \, n}\, d\omega .
\end{equation}
Asi obtenemos
\begin{equation}
  \begin{array}{l}
    \frac{1}{2 \, \pi }\int_{0}^{2 \, \pi}
    \left( \sum_{k=-\infty }^{\infty }
      2 \, \pi \, a_{k} \, \delta \left( \omega -\frac{2 \, \pi \, k}{N}\right) 
    \right) \, e^{j \, \omega \, n}\, d\omega = \\
\\
    \left( \sum_{k=-\infty }^{\infty }
      a_{k}\int_{0}^{2 \, \pi}\delta 
      \left( \omega -\frac{2 \, \pi \, k}{N}\right) e^{j \, \omega \, n}\, d\omega 
    \right).
  \end{array}
\end{equation}

\item Aplicando la propiedad del muestreo 
\begin{equation}
  \delta \left( \omega -\frac{2 \, \pi \, k}{N}\right) e^{j \, \omega \, n} = 
  \delta \left( \omega -\frac{2 \, \pi \, k}{N}\right) e^{j \frac{2 \, \pi \, k}{N} n},
\end{equation}
intercambiando la posici\'{o}n de las sumatorias y las integrales, y simplificando coeficientes
se obtiene
\begin{equation}
\begin{array}{l}
\sum_{k=-\infty }^{\infty }
a_{k}\int_{0}^{2 \, \pi }\delta \left( \omega -\frac{2 \, \pi \, k}{N}\right) e^{j \, \omega \, n}\, d\omega  = \\
\\
\sum_{k=-\infty}^{\infty }
a_{k} \, e^{j \, \frac{2 \, \pi \, k}{N} \, n}
\left( \int_{0}^{2 \, \pi}\delta \left(\omega -\frac{2 \, \pi \, k}{N}\right) \, d\omega \right).
\end{array}
\end{equation}

\item Evaluemos la integral de acuerdo al valor de $k$. Sabemos que se debe cumplir $\omega = \frac{2 \, \pi \, k}{N}$ en el intervalo de integración para que
la integral no sea nula
\begin{equation}
\int_{0}^{2 \, \pi}\delta \left( \omega -\frac{2 \, \pi \, k}{N}\right) \, d\omega
=
\left\{ 
\begin{array}{c}
1 \\ 
0
\end{array}
\begin{array}{c}
0\leq k<N, \\ 
\left( k<0\right) \vee \left(N \leq k\right).
\end{array}
\right.
\end{equation}
\item Evaluamos la sumatoria de integrales para obtener
\begin{equation}
\sum_{k=-\infty }^{\infty }
a_{k} \, e^{j\frac{2 \, \pi \, k}{N}n}
\left( \int_{0}^{2 \, \pi}\delta \left( \omega -\frac{2 \, \pi \, k}{N}\right) \, d\omega \right)=
\sum_{k=0}^{N-1}a_{k} \, e^{j \, \frac{2 \, \pi \, k}{N} \, n}.
\end{equation}

\item Obtenemos finalmente la s\'{i}ntesis
\begin{equation}
x\left[ n\right] =\frac{1}{2 \, \pi }\int_{0}^{2 \, \pi }X\left( e^{j \, \omega
}\right) e^{j \, \omega \, n}\, d\omega =\sum_{k=0}^{N-1}a_{k} \, e^{j\frac{2 \, \pi \, k}{N} \, n}.
\end{equation}

\end{enumerate}

\section{Propiedades de la Transformada de Fourier de tiempo discreto}

\subsection{Ecuaci\'{o}n de s\'{\i}ntesis}

\begin{equation}
x\left[ n\right] =
\frac{1}{2 \, \pi }\int_{0}^{2 \, \pi }X\left( e^{jk\omega}\right) e^{jk\omega \, n}\, d\omega.
\end{equation}

\subsection{Ecuaci\'{o}n de an\'{a}lisis}

\begin{equation}
X\left( j \, \omega \right) =
\sum_{n=-\infty }^{\infty }x\left[ n\right]e^{-j \, \omega \, n}.
\end{equation}

\subsection{Notaci\'{o}n}

\begin{equation}
x\left[ n\right] 
\begin{array}{l}
TF \\ 
\leftrightarrow
\end{array}
X\left( e^{j \, \omega }\right) .
\end{equation}

\subsection{Periodicidad}

\begin{equation}
X\left( e^{j\left( \omega +2 \, \pi \right) }\right) =
X\left( e^{j \, \omega}\right).
\end{equation}

\subsection{Linealidad}

\begin{equation}
x\left[ n\right] 
\begin{array}{l}
TF \\ 
\leftrightarrow
\end{array}
X\left( e^{j \, \omega }\right) ,
\end{equation}

\begin{equation}
y\left[ n\right] 
\begin{array}{l}
F \\ 
\leftrightarrow
\end{array}
Y\left( e^{j \, \omega }\right) ,
\end{equation}
\begin{equation}
a \, x\left[ n\right] +b \, y\left[ n\right] 
\begin{array}{l}
TF \\ 
\leftrightarrow
\end{array}
a \, X\left( e^{j \, \omega }\right) +b \, Y\left( e^{j \, \omega }\right) .
\end{equation}

\subsection{Desplazamiento en el tiempo}

\begin{equation}
x\left[ n\right] 
\begin{array}{l}
TF \\ 
\leftrightarrow
\end{array}
X\left( e^{j \, \omega }\right) ,
\end{equation}

\begin{equation}
x\left[ n-n_{0}\right] 
\begin{array}{l}
TF \\ 
\leftrightarrow
\end{array}
e^{-j \, \omega \, n_{0}}X\left( e^{j \, \omega }\right) ,
\end{equation}

\subsection{Desplazamiento en frecuencia}

\begin{equation}
x\left[ n\right] 
\begin{array}{l}
TF \\ 
\leftrightarrow
\end{array}
X\left( e^{j \, \omega }\right).
\end{equation}

\begin{equation}
e^{-j \, \omega_{0} \, n}x\left[ n\right] 
\begin{array}{l}
TF \\ 
\leftrightarrow
\end{array}
X\left( e^{j\left( \omega -\omega _{0}\right) }\right).
\end{equation}

\subsection{Conjugaci\'{o}n y simetr\'{\i}a del conjugado}

\begin{equation}
x\left[ n\right] 
\begin{array}{l}
TF \\ 
\leftrightarrow
\end{array}
X\left( e^{j \, \omega }\right) ,
\end{equation}

\begin{equation}
x^{ }\left[ n\right] =\bar{x}\left[ n\right] 
\begin{array}{l}
TF \\ 
\leftrightarrow
\end{array}
X^{ }\left( e^{-j \, \omega }\right) =\bar{X}\left( e^{-j \, \omega }\right) ,
\end{equation}

Si $x\left[ n\right] $ es real 
\begin{eqnarray}
&&Even\left( x\left[ n\right] \right) 
\begin{array}{l}
TF \\ 
\leftrightarrow
\end{array}
\Re\left[ X\left( e^{j \, \omega }\right) \right] , \\
&&Odd\left( x\left[ n\right] \right) 
\begin{array}{l}
TF \\ 
\leftrightarrow
\end{array}
\Im\left[ X\left( e^{j \, \omega }\right) \right] .
\end{eqnarray}

\subsection{Diferencia}

\begin{equation}
x\left[ n\right] 
\begin{array}{l}
TF \\ 
\leftrightarrow
\end{array}
X\left( e^{j \, \omega }\right) ,
\end{equation}

\begin{equation}
x\left[ n\right] -x\left[ n-1\right] 
\begin{array}{l}
TF \\ 
\leftrightarrow
\end{array}
\left( 1-e^{j \, \omega }\right) X\left( e^{j \, \omega }\right).
\end{equation}

\subsection{Suma}

\begin{equation}
x\left[ n\right] 
\begin{array}{l}
TF \\ 
\leftrightarrow
\end{array}
X\left( e^{j \, \omega }\right) ,
\end{equation}

\begin{equation}
\sum_{k=-\infty }^nx\left[ k\right] = 
\begin{array}{l}
F \\ 
\leftrightarrow
\end{array}
\frac{X\left( e^{j \, \omega }\right) }{\left( 1-e^{j \, \omega }\right) }+\pi
X\left( e^{j0}\right) \sum_{k=-\infty }^{\infty }\delta \left( \omega -2 \, \pi
k\right) .
\end{equation}

\subsection{Diferenciaci\'{o}n en frecuencia}

\begin{equation}
x\left[ n\right] 
\begin{array}{l}
TF \\ 
\leftrightarrow
\end{array}
X\left( e^{j \, \omega }\right) ,
\end{equation}

\begin{equation}
nx\left[ n\right] 
\begin{array}{l}
TF \\ 
\leftrightarrow
\end{array}
j\frac{dX\left( e^{j \, \omega }\right) }{\, d\omega }.
\end{equation}

\subsection{Relaci\'{o}n de Parseval}

\begin{equation}
\sum_{n=-\infty }^{\infty }\left| x\left[ n\right] \right| ^{2}=
\frac{1}{2 \, \pi }
\int_{0}^{2 \, \pi }\left| X\left( e^{j \, \omega }\right) \right| ^{2}\, d\omega.
\end{equation}

\subsection{Convoluci\'{o}n}

\begin{equation}
x\left[ n\right] 
\begin{array}{l}
TF \\ 
\leftrightarrow
\end{array}
X\left( e^{j \, \omega }\right) ,
\end{equation}

\begin{equation}
y\left[ n\right] 
\begin{array}{l}
TF \\ 
\leftrightarrow
\end{array}
Y\left( e^{j \, \omega }\right) ,
\end{equation}

\begin{equation}
h\left[ n\right] 
\begin{array}{l}
TF \\ 
\leftrightarrow
\end{array}
H\left( e^{j \, \omega }\right) ,
\end{equation}
\begin{equation}
y\left[ n\right] =h\left[ n\right]  x\left[ n\right] 
\begin{array}{l}
TF \\ 
\leftrightarrow
\end{array}
Y\left( e^{j \, \omega }\right) =H\left( e^{j \, \omega }\right) X\left( e^{j \, \omega
}\right) .
\end{equation}

\subsection{Multiplicaci\'{o}n}

\begin{equation}
x\left[ n\right] 
\begin{array}{l}
F \\ 
\leftrightarrow
\end{array}
X\left( e^{j \, \omega }\right) ,
\end{equation}

\begin{equation}
y\left[ n\right] 
\begin{array}{l}
F \\ 
\leftrightarrow
\end{array}
Y\left( e^{j \, \omega }\right) ,
\end{equation}

\begin{equation}
r\left[ n\right] 
\begin{array}{l}
F \\ 
\leftrightarrow
\end{array}
R\left( e^{j \, \omega }\right) ,
\end{equation}

\begin{equation}
r\left[ n\right] =x\left[ n\right] y\left[ n\right] 
\begin{array}{l}
TF \\ 
\leftrightarrow
\end{array}
R\left( e^{j \, \omega }\right) =
\frac{1}{2 \, \pi }
\int_{-\infty }^{\infty }
  X\left(e^{j\theta }\right) Y\left( e^{j\left( \omega -\theta \right) }\right)d\theta .
\end{equation}
